\section{INVERSE LIMITS OF MEASURES}

Throughout this section, I denotes a nonempty set, equipped with a preorder relation, denoted \(i \leqslant j\) , and directed for this relation. Recall (GT, I, §4, No.~4) that an inverse system of topological spaces indexed by I is a family \((\mathrm{T}_i, p_{ij})\) where \(\mathrm{T}_i\) is a topological space and \(p_{ij}\) is a continuous mapping of \(\mathrm{T}_j\) into \(\mathrm{T}_i\) for \(i \leqslant j\) , where \(p_{ii}\) is the identity mapping of \(\mathrm{T}_i\) , and where \(p_{ik} = p_{ij} \circ p_{jk}\) for \(i \leqslant j \leqslant k\) . Let T be a topological space and \((p_i)_{i \in I}\) a family of continuous mappings \(p_i: \mathrm{T} \to \mathrm{T}_i\) . The family \((p_i)_{i \in I}\) is said to be coherent if \(p_i = p_{ij} \circ p_j\) for \(i \leqslant j\) , and it is said to be separating if for distinct \(x, y\) in T there exists an \(i \in I\) such that \(p_i(x) \neq p_i(y)\) . When \(\mathrm{T} = \varprojlim \mathrm{T}_i\) and \(p_i\) is the canonical mapping of T into \(\mathrm{T}_i\) , the family \((p_i)_{i \in I}\) is coherent and separating.

